\documentclass[10pt,a4paper]{amsart}
\usepackage[margin=19mm]{geometry}
\usepackage{amsmath,amssymb,mathtools}
\usepackage[scr=rsfs,cal=euler]{mathalfa}
\usepackage{xfrac}
\usepackage[unicode,hidelinks,bookmarks=false]{hyperref}
\newcommand{\ie}{i.e.\ }
\newcommand{\cf}{cf.\ }
\newcommand{\resp}{respectively\ }
\newcommand{\Subsection}{\subsection}
\providecommand{\nobreakdash}{\nobreak\hbox{-}}
\setlength{\emergencystretch}{2em}
\multlinegap0pt
\usepackage{amssymb}
\usepackage{tcolorbox}
\usepackage{algorithm}\usepackage{algpseudocode}
\newtheorem{lemma}{Lemma}
\title{On the graph of the dimension function of the Lagrange and Markov spectra}
\author{Carlos Matheus, Carlos Gustavo Moreira, Polina Vytnova}
\date{Source: arXiv:2212.11371v2 (CC BY 4.0)}
\begin{document}
\maketitle

\noindent\emph{Source and licence.} On the graph of the dimension function of the Lagrange and Markov spectra. Version arXiv:2212.11371v2 (CC BY 4.0), distributed by arXiv under the Creative Commons Attribution 4.0 International licence, http://creativecommons.org/licenses/by/4.0/, as recorded in the arXiv licence metadata for this exact version and shown on the abstract page https://arxiv.org/abs/2212.11371v2. Full licence notice: \texttt{licenses/arxiv-2212.11371-CC-BY-4.0.txt}.\par\smallskip

\noindent\textit{Editorial scope.} Counted unit: the article's lemma l.plateaux (source0.tex lines 373--376). The article's complete original proof of it is delivered with it (source0.tex lines 378--380, one \texttt{proof} environment of 3 lines). No mathematics is written, repaired or re-derived: the statement and the proof are the article's own lines, in the article's own notation, and no retained line is substituted or re-flowed.  No further source-local context is needed: every cross-reference of every retained line resolves inside the two delivered spans.  Nothing was omitted from any retained span, so the package carries no omission record.  No retained line contains a citation, so the wrapper carries no bibliography and no bibliography entry.  No retained line contains a figure environment, an image-inclusion command or drawing code, so no image is delivered and the package declares no asset dependency.  Editorial formatting only, in addition: an AMS \texttt{amsart} preamble that restates the article's own declarations and macros used by the retained lines, namely the environments \texttt{lemma} on separate counters, together with the standard packages those lines need.  The retained environments print in this document as Lemma 1 in order of appearance, whereas the article prints the same items with the numbering of its own full document.  The complete native source of the article as distributed in the arXiv e-print for this exact version is bundled unchanged as \texttt{sources/135211/source0.tex}.
\begin{lemma}
    \label{l.plateaux} 
    Let $\beta$ be a finite word and $\theta$ be a symmetric finite word of even size on $\{1,2\}$. Assume that, under some conditions (e.g., after forbidding a finite list of finite strings), the Markov value of an infinite word of the type $\gamma=\omega_2^t \beta^t 2\theta 2\beta \omega_1$ is attained at one of the $2$ next to $\theta$, where $\omega_1, \omega_2$ are infinite words (on $\{1,2\}$). Then this Markov value is minimum when $\omega_1=\omega_2=:\omega$ and $[0;\beta,\omega]$ is minimum (under these conditions).
\end{lemma} 

\begin{proof} Let $x=[0;\beta,\omega_1]$ and $y=[0;\beta,\omega_2]$. If the positions with the $2$'s after and before $\theta$ are, respectively, $0$ and $2r+1$ (where $2r=|\theta|$), then $\lambda_0(\gamma)=2+y+[0;\theta,2+x]=:h(x,y)$ and $\lambda_{2r+1}(\gamma)=2+x+[0;\theta,2+y]=h(y,x)$. Assume without loss of generality (by symmetry) that the Markov value of $\gamma$ is attained at the position $0$. Since $h(x,y)$ is increasing in $y$ and decreasing in $x$, then we should have $y\ge x$ (indeed if $y<x$ then $\lambda_0(\gamma)=h(x,y)<h(x,x)<h(y,x)=\lambda_{2r+1}(\gamma)$, a contradiction). 
Consequently, in this case we have $h(x,y)\ge h(y,y)$, so the minimum value in this case is attained when $x=y$, and is equal to $h(y,y)=2+y+[0;\theta,2+y]$. Since $f(y)=[0;\theta,2+y]$ is a contraction, $h(y,y)$ is increasing, so the minimum is attained when $y=[0;\beta,\omega_2]=[0;\beta,\omega]$ is minimum. 
\end{proof} 

\end{document}
